\section{Efficient probes and the prediction curve}
\label{sec:prediction}

Fix $s=1/8$ and take one of the two subsequences supplied by
Proposition~\ref{prop:critical-sequence}.  We continue to suppress the
counterexample index.  Thus $\lambda\to1$, $\eta$ satisfies the logarithmic
Sobolev inequality with constant at most $1/\lambda$, and, with
$A=A_\eta$, $w$ belongs coefficientwise to $\Dom(A)$ and satisfies
\begin{equation}
 I=\E_\eta|w|^2>0,
 \qquad \|(A-1)w\|_{L^2(\eta)}+|\E_\eta w|=o(\sqrt I).
 \label{eq:prediction-input}
\end{equation}
In the strict case, $\eta=f^2\mu_s$, $w=D\log f^2$, and
$2H(\eta\mid\mu_s)=I/\lambda$.  In the gap case,
$\eta=\mu_s$, $w=D\varphi$, and
$A_{\mu_s}\varphi=\lambda\varphi$, where
$\E_{\mu_s}\varphi=0$ and $I=\lambda\E_{\mu_s}\varphi^2$.
All limits in this section are along the selected subsequence, with the
displayed auxiliary parameters fixed unless a different order is specified.

Our objective is to determine how much of $w(Y_0)$ can be predicted from
$Y_0+\sqrt rG$.  The strict case first requires an entropy estimate for an
additional Gaussian observation of $w$.  We then prove an entropy-loss
comparison for arbitrary smooth alternatives to a log-concave law.  These
two estimates determine the prediction curve; linearization of the same
comparison handles the gap case.

\begin{lemma}[Efficient Gaussian probes]
\label{lem:efficient-probes}
In the strict case, let $Y_0\sim\eta$ and let $G'$ be an independent
standard Gaussian vector in the same Euclidean space.  For $T\ge0$ define
\[
 L_T=\sqrt T\,w(Y_0)+G',\qquad
 d(T)=H(\Law(L_T)\mid\gamma),
\]
where $\gamma$ is standard Gaussian measure.  For every finite $M$,
\begin{equation}
 \sup_{0\le T\le M}d(T)=o(I).
 \label{eq:efficient-probes}
\end{equation}
\end{lemma}

\begin{proof}
Run the stationary diffusion with law $\eta$, using Brownian noise
independent of $G'$, and write $P_u=e^{-uA}$.
Proposition~\ref{prop:critical-sequence} puts each coordinate in
$\Dom(A)$ and gives the identity
\begin{equation}
 -Ay=-y/s+w+b_s/s.
 \label{eq:prediction-coordinate-generator}
\end{equation}
In particular, $\E_\eta|Y_0|^2<\infty$. The logarithmic Sobolev
inequality for $\eta$ then gives
\[
 \E_\eta e^{v\cdot(Y_0-\E Y_0)}
 \le e^{|v|^2/(2\lambda)}.
\]
Integrating this inequality over a Gaussian $v$ proves
$\E_\eta e^{c'|Y_0|^2}<\infty$ for a sufficiently small $c'>0$.
Here the choice of $c'$ may depend on the individual law and dimension.
This exponential-square moment will also permit application of the
entropy-loss comparison below.

For a fixed $T\le M$, set
\[
 J_u=\MI(L_T;Y_u),\qquad
 F_u=\E\,\Ical(\Law(Y_u\mid L_T)\mid\eta),\qquad
 Q_u=\E|\E[Y_u\mid L_T]|^2.
\]
The Gaussian observation identity gives
$J_0=TI/2-d(T)$.  Entropy dissipation, including its conditional version,
gives $-dJ_u\ge F_u\,du$, and in particular
$\int_0^1F_u\,du\le J_0$.  If $\sigma_{u,l}=\Law(Y_u\mid L_T=l)$,
then the entropy chain rule and the strict-case normalization give
\begin{equation}
 \E H(\sigma_{u,L_T}\mid\mu_s)=J_u+\frac{I}{2\lambda},
 \qquad
 \E\Ical(\sigma_{u,L_T}\mid\mu_s)=F_u+I.
 \label{eq:probe-reference-decomposition}
\end{equation}
For the second equality, write the conditional relative score with
respect to $\mu_s$ as the sum of $w$ and the conditional relative score
with respect to $\eta$.  The latter scores have conditional average zero
at each $Y_u$, since the conditional likelihoods average to one.  This
cancels the cross term.  At the almost every time where $F_u<\infty$,
the assertion also follows weakly by differentiating that marginal
identity: Cauchy--Schwarz gives locally integrable averaged derivatives.

Apply Lemma~\ref{lem:translation-gain} to $\sigma_{u,l}$ and average
over $l$.  Since $s^{-2}-\lambda/s$ stays positive, Equations
\eqref{eq:probe-reference-decomposition} give, on $[0,1]$,
\begin{equation}
 -dJ_u\ge(2\lambda J_u+c_sQ_u)\,du-d\mathcal E_u,
 \qquad \mathcal E([0,1])=o_M(I),
 \label{eq:probe-decay}
\end{equation}
where $c_s>0$ is fixed and $\mathcal E$ is a nonnegative error measure.
Indeed the error in that lemma is bounded by a number tending to zero
times $F_u+I+J_u+I/(2\lambda)$, whose integral is $O_M(I)$.
The measure formulation avoids imposing differentiability on $J_u$.

The spectral theorem and Equation~\eqref{eq:prediction-input} imply
\begin{equation}
 \sup_{0\le u\le1}\|P_uw-e^{-u}w\|_2
 \le \|(A-1)w\|_2=o(\sqrt I).
 \label{eq:probe-semigroup}
\end{equation}
Reversibility gives
$\E[L_T\mid Y_u]=\sqrt T P_uw(Y_u)$.  Relative entropy with respect
to a standard Gaussian is at least half the square of the mean, by
testing its variational formula with linear functions.  Consequently
\[
 J_u+d(T)=\E H(\Law(L_T\mid Y_u)\mid\gamma)
 \ge\frac T2\|P_uw\|_2^2
 =\frac{TI}{2}e^{-2u}+o_M(I).
\]
Integrating Equation~\eqref{eq:probe-decay} with factor $e^{2\lambda u}$
and comparing at $u=1$ yields
\begin{equation}
 \int_0^1Q_u\,du\le C_s d(T)+o_M(I).
 \label{eq:probe-position-control}
\end{equation}
Here the difference between $e^{-2\lambda}$ and $e^{-2}$ contributes
only $o_M(I)$.

It remains to turn control of the predicted position into control of the
predicted score.  The pushforward of $\mu_s$ by $b_s$ has the centered
linear moment bound
\[
 \log\E_{\mu_s}e^{v\cdot b_s}\le Ca^2|v|^2,
\]
by conditional Jensen and the corresponding bound for the original
signal $X$.  The entropy variational formula, optimized over $v$, thus
gives
\[
 \E|\E[b_s(Y_v)\mid L_T]|^2
 \le Ca^2\bigl(J_v+I/(2\lambda)\bigr)=o_M(I).
\]
Condition Equation~\eqref{eq:prediction-coordinate-generator} on $L_T$
and use variation of constants.  By
Equation~\eqref{eq:probe-semigroup} and the contraction property of
conditional expectation,
\begin{equation}
 \E[Y_u\mid L_T]
 =e^{-u/s}\E[Y_0\mid L_T]
 +\frac{e^{-u}-e^{-u/s}}{1/s-1}\E[w(Y_0)\mid L_T]
 +o_{L^2,M}(\sqrt I),
 \label{eq:probe-coordinate-solution}
\end{equation}
uniformly for $0\le u\le1$.  The two scalar functions multiplying the
conditional means are linearly independent on $[0,1]$: the first is
nonzero at zero, whereas the second vanishes there and is nonzero for
$u>0$.  Their $2\times2$ Gram matrix is therefore positive definite.
Applying its smallest-eigenvalue bound to each coordinate and then
averaging, Equations~\eqref{eq:probe-position-control} and
\eqref{eq:probe-coordinate-solution} show that
\[
 \E|\E[w(Y_0)\mid L_T]|^2\le C_s d(T)+o_M(I).
\]
By Lemma~\ref{lem:gaussian-signals}, the left side is $2d'(T)$ for
almost every $T$.  The error is uniform for $T\in[0,M]$, and $d(0)=0$.
Integrating this differential inequality proves
Equation~\eqref{eq:efficient-probes}.  No lower bound on $I$ is used.
\end{proof}

The next comparison relates the entropy lost under heat flow to the
conditional mean of the initial relative score.  The conditional mean
need not equal the relative score at positive heat time, so entropy
dissipation alone does not give this comparison.

The choice of observation directions is related to function-preserving
stochastic localization. Eldan--Koehler--Zeitouni
\cite[Section~2, Lemma~2]{EldanKoehlerZeitouni2022} use controls
approximating projection orthogonal to the covariance of the signal
with a selected function. Chen and Eldan
\cite[Sections~2.4.2 and~3.2.3, Equation~(27)]{ChenEldan2022}
express the entropy drift through the difference of means under a
measure and its normalized likelihood tilt, in a continuous-time
setup with positive definite controls. Below we prove the cancellation
and its error directly for finite Gaussian observations with singular
projections, and use independent copies to control the discarded
directions.

\begin{lemma}[Entropy loss under a Gaussian observation]
\label{lem:entropy-loss}
Let $p$ and $q$ be smooth strictly positive probability densities on
$\R^d$, with $q$ log-concave.  Suppose
\[
 H_0=H(p\mid q)<\infty,\qquad
 w=D\log(p/q),\qquad I=\E_p|w|^2<\infty,
 \qquad \E_p e^{c|Y|^2}<\infty
\]
for some $c>0$, where $Y\sim p$. Let $G$ be a standard Gaussian
vector independent of $Y$. For $r>0$, let
$U_r=Y+\sqrt rG$, let $p_r$ and $q_r$ be the Gaussian convolutions
of $p$ and $q$, and put
$\bar w_r(U_r)=\E_p[w(Y)\mid U_r]$ and $H_r=H(p_r\mid q_r)$.
Then, for every $z>0$,
\begin{equation}
 2(H_0-H_z)\le\int_0^z\E_p|\bar w_r(U_r)|^2\,dr.
 \label{eq:prediction-entropy-loss}
\end{equation}
\end{lemma}

\begin{proof}
The proof adds weak observations that increase the curvature of the
base posterior while losing asymptotically no relative entropy.
At each step we remove the direction in which the two posterior means
differ. The retained output means then agree, and we will show that
its relative entropy, averaged over the retained history, is
$O(\delta^2)$ for observation strength $\delta$.
Independent copies make the cost of removing these directions vanish
after division by the number of copies. That last comparison concerns
the score residuals after full observations, whose conditional
independence will be essential.

We keep the one-copy problem fixed. The three successive limits will
be refinement of a finite observation mesh, increase of the number of
independent copies, and increase of the total additional precision.
Set $b=1/z$ and take $N$ independent
copies.  Write $Y=(Y_1,\ldots,Y_N)$ and
$W=(w(Y_1),\ldots,w(Y_N))$; the two possible priors are $p^{\otimes N}$
and $q^{\otimes N}$.  Expectations in the proof use $p^{\otimes N}$.

\emph{Additional observations and their entropy cost.}
First observe $Z_0=\sqrt bY+G_0$.  The chain rule says that the
remaining posterior relative entropy, averaged over this observation,
is $N(H_0-H_z)$.  Fix $u>0$ and an integer $m$, and put $\delta=u/m$.
For $i=0,\ldots,m-1$, let $\mathcal H_i$ be the retained history before
step $i$.  Compute the difference of the two posterior means given
this history and let $P_i$ be the orthogonal projection onto its span;
set $P_i=0$ when the difference vanishes.  Retain only
\[
 (1-P_i)Z_{i+1},\qquad Z_{i+1}=\sqrt\delta Y+G_{i+1},
\]
where all Gaussian noises are independent.  The same measurable rule
is used under both priors.  One may choose versions of the posterior
means by their integral formulas; the positive densities make the rule
well defined on all histories of positive probability.

For each realized retained history the likelihood, as a function of
$Y$, is a Gaussian quadratic factor of precision
\begin{equation}
 B=b\Id+\delta\sum_{i=0}^{m-1}(1-P_i)\succeq b\Id.
 \label{eq:entropy-loss-precision}
\end{equation}
There is no extra likelihood factor from the choices of $P_i$: each
choice was already determined by the retained past.  The factor is the
same for both priors.  Thus the posterior base has curvature at least
$B$, and the posterior relative score is still $W$.

We next bound the relative entropy lost in the additional observations.
Condition on a retained history at one step and put $Q=1-P_i$.
The projected posterior means agree.  After subtracting their common
mean, let $X_c$ denote either centered projected signal in
$\operatorname{ran}Q$, whose dimension is at most $Nd$.  The base
posterior has curvature at least $b\Id$.  Its centered linear moment
bound and integration over a Gaussian vector give
\[
 \E_q[e^{\kappa|X_c|^2}\mid\mathcal H_i]
 \le (1-2\kappa/b)^{-Nd/2},\qquad 0<2\kappa<b.
\]
For the alternative, write $a_i=\E_p[Y\mid\mathcal H_i]$.  For any
sufficiently small $\kappa>0$, conditional Jensen gives
\begin{align*}
 \E_p[e^{\kappa|Q(Y-a_i)|^2}\mid\mathcal H_i]
 &\le e^{2\kappa|a_i|^2}
       \E_p[e^{2\kappa|Y|^2}\mid\mathcal H_i]\\
 &\le\bigl(\E_p[e^{2\kappa|Y|^2}\mid\mathcal H_i]\bigr)^2.
\end{align*}
Squaring and applying conditional Jensen once more shows that its
square, averaged over $\mathcal H_i$, is bounded by
$\E_p e^{8\kappa|Y|^2}<\infty$.  These bounds hold for every history
sigma-field, so their constants do not depend on the mesh or on the
adaptive choices.

Let $h_p,h_q$ be the output densities of
$\sqrt\delta X_c+G$ in $\operatorname{ran}Q$ and let $\phi$ be standard
Gaussian density there.  Centering and Jensen give the pointwise bound
\[
 \frac{h_q(x)}{\phi(x)}
 =\E_q e^{\sqrt\delta x\cdot X_c-\delta|X_c|^2/2}
 \ge e^{-\delta Nd/(2b)}.
\]
For either conditional signal, using a conditionally independent copy
$X_c'$ and integrating the Gaussian kernel gives
\[
 \left\|\frac h\phi-1\right\|_{L^2(\phi)}^2
 =\E[e^{\delta X_c\cdot X_c'}\mid\mathcal H_i]-1.
\]
The linear term vanishes because $X_c$ is centered.  The bound
$|e^t-1-t|\le t^2e^{|t|}/2$, followed by
$|X_c\cdot X_c'|\le(|X_c|^2+|X_c'|^2)/2$, bounds the last expression
by $C\delta^2$ times the square of a small conditional
exponential-square moment.  The preceding moment bounds therefore give
\[
 \E H(h_p\mid h_q)
 \le\E\chi^2(h_p\mid h_q)
 \le2e^{\delta Nd/(2b)}\E\left(
 \left\|h_p/\phi-1\right\|_2^2+
 \left\|h_q/\phi-1\right\|_2^2\right)
 \le C_{N,b,p}\delta^2.
\]
This statement includes the zero-dimensional observed subspace, for
which the loss is zero.  Summing the $m$ steps loses
$O_{N,b,p}(u\delta)$ entropy.  Apply the anisotropic curvature
inequality, Lemma~\ref{lem:curvature}, to the final posterior pair.
For fixed $N$ and $u$ we obtain
\begin{equation}
 2N(H_0-H_z)-o_{\delta\downarrow0}(1)
 \le\E W^*B^{-1}W.
 \label{eq:entropy-loss-posterior-lsi}
\end{equation}

\emph{Comparison with full observations.}
We now estimate the effect of the rank-one exclusions. Work in the
Hilbert space of square-integrable vector fields on the signal and all
the full observations $Z_0,\ldots,Z_m$.  In particular,
$\|W\|_2^2=NI$.  Completing the square pointwise gives
\begin{equation}
 b\|W\|_2^2-b^2\E W^*B^{-1}W
 =\min_v\left\{b\|W-v\|_2^2+
              \delta\sum_{i=0}^{m-1}\|(1-P_i)v\|_2^2\right\}.
 \label{eq:entropy-loss-variational}
\end{equation}
Its minimizer is $v_*=bB^{-1}W$, so $\|v_*\|_2\le\|W\|_2$.
Let $\mathcal F_i=\sigma(Z_0,\ldots,Z_i)$ and let $e_i$ be
conditional expectation onto $\mathcal F_i$.  The sigma-fields are
nested, and $P_i$ is $\mathcal F_i$-measurable because the retained
history is a function of this full history. We compare the original
minimum with one that penalizes the part of $v$ left unpredicted by
the full observations: replace $1-P_i$ by $1-e_i$, and denote the
resulting minimum by $M_0$. The nested projections make this new
minimum explicitly computable in terms of Gaussian prediction.
We first find its minimizer, then use its residuals in a dual test
to bound the original minimum below in terms of $M_0$.
Its minimizer is
\[
 v_0=b\left(b+\delta\sum_{i=0}^{m-1}(1-e_i)\right)^{-1}W.
\]
Here is its explicit form.  Set $\alpha_k=b/(b+k\delta)$ for
$0\le k\le m$.  On $\operatorname{ran}e_0$ the multiplier is one;
on $\operatorname{ran}(e_k-e_{k-1})$ it is $\alpha_k$; on
$\operatorname{ran}(1-e_{m-1})$ it is $\alpha_m$.  Consequently
\begin{equation}
 v_0=\alpha_mW+
       \sum_{k=0}^{m-1}(\alpha_k-\alpha_{k+1})e_kW.
 \label{eq:entropy-loss-minimizer}
\end{equation}
All coefficients are nonnegative and sum to one.  Put
$g_i=(1-e_i)v_0$.  Equation~\eqref{eq:entropy-loss-minimizer} gives
\begin{equation}
 g_i=\alpha_m(W-e_iW)+
 \sum_{k=i+1}^{m-1}(\alpha_k-\alpha_{k+1})(e_kW-e_iW).
 \label{eq:entropy-loss-residual}
\end{equation}
The weights in this formula sum to $\alpha_{i+1}\le1$.

We need a bound uniform over every subdivision.  The retained histories need not factor across copies.  We use
independence only for the full observation arrays: those belonging to
different copies are independent under $p^{\otimes N}$.  Conditioning
on $\mathcal F_i$ conditions each such array on its own past.  Thus
$e_kW_j$ depends only on copy $j$'s observations
through step $k$.  Conditional on $\mathcal F_i$, the blocks $g_{i,j}$
are independent and centered, even though they involve future full
observations.  To see this directly, factor the joint law into the
independent copy arrays and condition each factor on its own past;
the resulting conditional kernel remains a product.  By conditional
Jensen applied to Equation~\eqref{eq:entropy-loss-residual},
\begin{equation}
 Z_{i,j}:=\tr\Cov(g_{i,j}\mid\mathcal F_i)
 \le\E[|W_j-e_iW_j|^2\mid\mathcal F_i]
 \le\E[|W_j|^2\mid\mathcal F_i].
 \label{eq:entropy-loss-block-bound}
\end{equation}
For example, every term with $k>i$ is bounded in conditional squared
mean by $\E[|W_j-e_iW_j|^2\mid\mathcal F_i]$ by the tower property;
the nonnegative weights of total at most one preserve this bound.

The right sides in Equation~\eqref{eq:entropy-loss-block-bound} are
uniformly integrable over all these sigma-fields.  Explicitly, if
$Z\ge0$ is integrable and $M=\E[Z\mid\mathcal F]$, then
\begin{equation}
 \E[M\1_{\{M>K\}}]\le2\E[Z\1_{\{Z>K/2\}}].
 \label{eq:entropy-loss-ui}
\end{equation}
Indeed $M\le K/2+\E[Z\1_{\{Z>K/2\}}\mid\mathcal F]$, and on
$\{M>K\}$ the last conditional expectation is at least $M/2$.
The conditional covariance of $g_i$ is block diagonal.  Since $P_i$
is measurable in $\mathcal F_i$ and has rank at most one,
\[
 \E|P_ig_i|^2\le\E\max_{1\le j\le N}Z_{i,j}.
\]
The maximum is at most $K+\sum_jZ_{i,j}\1_{\{Z_{i,j}>K\}}$.
Equations~\eqref{eq:entropy-loss-block-bound} and
\eqref{eq:entropy-loss-ui} show that, uniformly in $i,m$, and $\delta$,
\begin{equation}
 \E|P_ig_i|^2\le Nq_N,\qquad
 q_N:=\inf_{K>0}\left\{\frac KN+
           2\E_p[|w(Y)|^2\1_{\{|w(Y)|^2>K/2\}}]\right\}\longrightarrow0.
 \label{eq:entropy-loss-rank-error}
\end{equation}
No moment of the score beyond its second moment is required.  In
particular this bound permits $P_i$ to depend adaptively on all copies.

For completeness, the comparison of the two minima follows from a
dual test.  For every $v$,
\[
 \|(1-P_i)v\|_2^2
 \ge2\langle v,(1-P_i)g_i\rangle-\|g_i\|_2^2.
\]
The normal equation for $v_0$ is
$b(W-v_0)=\delta\sum_i g_i$.  Therefore the minimum over $v$ of
\[
 b\|W-v\|_2^2+2\delta\sum_i\langle v,g_i\rangle
                  -\delta\sum_i\|g_i\|_2^2
\]
is attained at $v_0$ and equals $M_0$.  Use the displayed linear bound
at $v_*$ in Equation~\eqref{eq:entropy-loss-variational}.  By
Equation~\eqref{eq:entropy-loss-rank-error}, the true minimum is at least
\begin{equation}
 M_0-2\delta\sum_i\|v_*\|_2\|P_ig_i\|_2
 \ge M_0-2uN\sqrt{Iq_N}.
 \label{eq:entropy-loss-dual-error}
\end{equation}

\emph{The prediction integral.}
It remains to compute $M_0$ and take the limits. The same orthogonal
decomposition into nested projection increments gives
\begin{equation}
 \frac{M_0}{N}
 =\sum_{i=0}^{m-1}
 \frac{b^2\delta}{(b+i\delta)(b+(i+1)\delta)}
       \left(I-\|e_iW_1\|_2^2\right).
 \label{eq:entropy-loss-discrete-minimum}
\end{equation}
Indeed a component first seen after $j$ increments has penalty
$bj\delta/(b+j\delta)$, and
\[
 \frac{bj\delta}{b+j\delta}
 =\sum_{i=0}^{j-1}
   \frac{b^2\delta}{(b+i\delta)(b+(i+1)\delta)}.
\]
Full observations through step $i$ have the sufficient Gaussian
statistic of precision $b+i\delta$.  Completing the likelihood square
proves sufficiency directly.  Hence
\[
 \|e_iW_1\|_2^2
 =\E_p|\bar w_{1/(b+i\delta)}(U_{1/(b+i\delta)})|^2.
\]
This squared prediction norm is nondecreasing with precision, since a
weaker observation can be obtained by adding independent Gaussian
noise.  It is bounded by $I$, so Equation
\eqref{eq:entropy-loss-discrete-minimum} converges as $m\to\infty$ to
\[
 b^2\int_0^u
 \frac{I-\E_p|\bar w_{1/(b+t)}(U_{1/(b+t)})|^2}{(b+t)^2}\,dt.
\]
First take this mesh limit at fixed $N,u$.  It removes the entropy
error in Equation~\eqref{eq:entropy-loss-posterior-lsi}, while
Equation~\eqref{eq:entropy-loss-dual-error} remains uniform in the
mesh.  Divide by $N$, then let $N\to\infty$, using $q_N\to0$.
Equations~\eqref{eq:entropy-loss-posterior-lsi} and
\eqref{eq:entropy-loss-variational} give
\[
 2(H_0-H_z)
 \le\frac Ib-\int_0^u
 \frac{I-\E_p|\bar w_{1/(b+t)}(U_{1/(b+t)})|^2}{(b+t)^2}\,dt
 =\frac{I}{b+u}
   +\int_{1/(b+u)}^{1/b}\E_p|\bar w_r(U_r)|^2\,dr.
\]
Finally let $u\to\infty$.  Since $1/b=z$ and the integrand is bounded
by $I$, this proves Equation~\eqref{eq:prediction-entropy-loss} in the
claimed order of limits.
\end{proof}

We return to the critical sequence.  For $Y_0\sim\eta$ and independent
Gaussian noise, define
\[
 U_r=Y_0+\sqrt rG,\qquad
 \bar w_r(U_r)=\E[w(Y_0)\mid U_r],\qquad
 m(r)=I^{-1}\E|\bar w_r(U_r)|^2.
\]
In the strict case we also use the label $L_T$ of
Lemma~\ref{lem:efficient-probes}, constructed with fresh Gaussian noise
independent of $(Y_0,G)$. Thus $L_T$ and $U_r$ are conditionally
independent given $Y_0$. Set
\[
 p_T(r)=I^{-1}\E|\E[w(Y_0)\mid U_r,L_T]|^2,
 \qquad p_0(r)=m(r).
\]

\begin{proposition}[The prediction curve]
\label{prop:prediction}
For the strict and gap critical sequences, for every fixed
$0<r<1/4$,
\begin{equation}
 m(r)\longrightarrow e^{-r}.
 \label{eq:prediction-curve}
\end{equation}
In the strict case, for every fixed finite $T\ge0$ one also has
$p_T(r)\to e^{-r}$.  In the gap case the eigenfunction may be rescaled
so that $I\to0$, without changing Equation~\eqref{eq:prediction-curve}
or Equation~\eqref{eq:prediction-input}.
\end{proposition}

\begin{proof}
For $0\le r<1/4$ put
\[
 \theta(r)=\exp\left(-\int_0^r\frac{dt}{R_{s+t}}\right).
\]
Proposition~\ref{prop:regularized-constants} gives
$\theta(r)=e^{-r}+o(1)$, uniformly on compact intervals under
consideration.  The common heat flow dissipates relative entropy at
one half the relative Fisher information.  Thus its entropy relative
to $\mu_{s+r}$ is at most $\theta(r)$ times its initial entropy,
by the logarithmic Sobolev inequality at the contemporaneous base law.
This assertion applies to conditional laws after averaging as well.

In the strict case write
$H_r=H(\Law(U_r)\mid\mu_{s+r})$ and
$H_0=I/(2\lambda)$.  The entropy of $U_r$ given $L_T$, relative to
$\mu_{s+r}$ and averaged over the label, is at most
$\theta(r)(H_0+TI/2)$, since initially it is $H_0+\MI(Y_0;L_T)$.
The chain rule gives
\begin{align*}
 \frac I2\int_0^T p_v(r)\,dv
 &=\E H(\Law(L_T\mid U_r)\mid\gamma)\\
 &=\E H(\Law(U_r\mid L_T)\mid\mu_{s+r})+d(T)-H_r\\
 &\le\frac I2\theta(r)(1/\lambda+T)+d(T).
\end{align*}
The first equality is Lemma~\ref{lem:gaussian-signals}, conditionally
on $U_r$.  Fix $v\ge0$ and then $T>v$.  Since $p_t(r)$ is
nondecreasing in signal strength $t$, Lemma
\ref{lem:efficient-probes} yields
\[
 (T-v)\limsup p_v(r)\le e^{-r}(1+T).
\]
The sequence limit is taken with $T$ fixed.  Allowing this fixed $T$
to be arbitrarily large proves
\begin{equation}
 \limsup p_v(r)\le e^{-r}.
 \label{eq:prediction-upper}
\end{equation}

Lemma~\ref{lem:entropy-loss} applies to $p=\eta$ and $q=\mu_s$;
the exponential-square moment was proved in the proof of
Lemma~\ref{lem:efficient-probes}.  Since $H_z\le\theta(z)H_0$, it gives
\[
 \int_0^z m(r)\,dr\ge\frac{1-\theta(z)}{\lambda}
 \qquad(0<z<1/4).
\]
The reverse limiting inequality follows from
Equation~\eqref{eq:prediction-upper}, $0\le m\le1$, and reverse
Fatou.  Hence
\begin{equation}
 \int_0^z m(r)\,dr\longrightarrow1-e^{-z}.
 \label{eq:prediction-integral}
\end{equation}
The function $m$ is nonincreasing in $r$.  For $0<h<r$ with
$r+h<1/4$, its value at $r$ is bounded below and above by its
averages on $[r,r+h]$ and $[r-h,r]$, respectively.  Equation
\eqref{eq:prediction-integral} determines the limits of both averages;
letting $h\downarrow0$ gives Equation~\eqref{eq:prediction-curve}.
Finally $m(r)\le p_T(r)$ and Equation~\eqref{eq:prediction-upper}
give the assertion for every fixed $T$.

We now prove the gap assertion without using additional probes.
First linearize Lemma~\ref{lem:entropy-loss}.  Let $q$ be a bounded
smooth centered test under $\mu_s$, with bounded gradient, and let
\[
 d\eta_\varepsilon=(1+\varepsilon q)\,d\mu_s,
 \qquad |\varepsilon|\|q\|_\infty<1.
\]
This alternative has a finite exponential-square moment.  Write
$\bar q_r=\E_{\mu_s}[q(Y_0)\mid U_r]$ and
$\overline{Dq}_r=\E_{\mu_s}[Dq(Y_0)\mid U_r]$.
The output density ratio is $1+\varepsilon\bar q_r$, while the
predicted initial relative score is exactly
\[
 \E_{\eta_\varepsilon}[D\log(1+\varepsilon q)(Y_0)\mid U_r]
 =\frac{\varepsilon\overline{Dq}_r}{1+\varepsilon\bar q_r}.
\]
Consequently its squared norm averaged under the alternative output
is
$\varepsilon^2\E_{\mu_{s+r}}[|\overline{Dq}_r|^2/
(1+\varepsilon\bar q_r)]$.  The entropy expansion
$H((1+\varepsilon q)\mu_s\mid\mu_s)
=\varepsilon^2\|q\|_2^2/2+o(\varepsilon^2)$, and its output version,
now give
\begin{equation}
 \|q\|_2^2-\|\bar q_z\|_2^2
 \le\int_0^z\|\overline{Dq}_r\|_2^2\,dr.
 \label{eq:prediction-linearized-loss}
\end{equation}
Indeed the denominator is bounded away from zero uniformly in $r$,
and conditional Jensen bounds the integrand by a constant times
$\|Dq\|_2^2$, justifying the limit.  Approximation in the form norm
extends Equation~\eqref{eq:prediction-linearized-loss} to $q=\varphi$:
centered compact smooth approximants have the required boundedness,
and all conditional expectation maps are $L^2$ contractions.  To retain
compact support while centering, subtract the approximant's mean times
one fixed compact smooth function of mean one; this correction tends
to zero in the form norm.

For clarity, we also give the heat variance identity used here.  For
a centered scalar or finite-dimensional vector field $g\in L^2(\mu_s)$,
let $g_r=\E[g(Y_0)\mid U_r]$ under the base law.  Then
\begin{equation}
 \frac d{dr}\E_{\mu_{s+r}}|g_r|^2
 =-\E_{\mu_{s+r}}|Dg_r|^2.
 \label{eq:prediction-heat-variance}
\end{equation}
For bounded $g$, both the numerator density $g_r\mu_{s+r}$ and the
denominator density $\mu_{s+r}$ solve the heat equation, and the
quotient calculation with the square function gives
$(\partial_r-\Delta/2)(\mu_{s+r}|g_r|^2)
=-\mu_{s+r}|Dg_r|^2$.  Integrate against expanding cutoffs, whose
Laplacians tend uniformly to zero.  The bound
$|g_r|\le\|g\|_\infty$ removes the cutoff terms and gives the
integrated identity on positive time intervals.  Its initial endpoint
follows from $\|g_r(U_r)-g(Y_0)\|_2\to0$: for a bounded continuous
test use $g(U_r)$ as a predictor and dominated convergence under the
coupling $U_r=Y_0+\sqrt rG$, and then approximate in $L^2(\mu_s)$
and use conditional contraction.  Approximate a general $L^2$ field by bounded
ones.  Applying the same identity to their differences shows that
their gradients are Cauchy in the space-time energy norm, while
conditional contraction gives convergence of both endpoint norms.
This proves the integrated version of Equation
\eqref{eq:prediction-heat-variance} in full generality.

Poincar\'e at $\mu_{s+r}$ and preservation of the mean therefore imply
\[
 \E_{\mu_{s+r}}|g_r|^2\le\theta(r)\E_{\mu_s}|g|^2.
\]
Apply this first to $\varphi$, and then to $w-\E w$.
The latter application and Equation~\eqref{eq:prediction-input} give
$m(r)\le\theta(r)+o(1)$.  The former and
$\|\varphi\|_2^2=I/\lambda$, inserted into Equation
\eqref{eq:prediction-linearized-loss}, give
$\int_0^z m(r)\,dr\ge(1-\theta(z))/\lambda$.
The same integral squeeze and monotonicity argument prove
Equation~\eqref{eq:prediction-curve}.  Multiplying $\varphi$ by any
positive scalar multiplies every squared quantity by its square, so
we may finally choose these scalars to arrange $I\to0$.
\end{proof}

The hierarchy construction will use the following consequences of the
prediction curve.  The estimates distinguish conditional covariance
tails, which are needed only in the strict case, from bounds on
deterministic subspaces, which hold in both cases.

\begin{lemma}[Posterior covariance and outlier subspaces]
\label{lem:posterior-outliers}
Use the critical sequence and notation of
Proposition~\ref{prop:prediction}, with $I\to0$ in the gap case.
The following assertions hold.
\begin{enumerate}[label=\textup{(\roman*)}]
\item For every fixed $c<\infty$, uniformly over deterministic
orthogonal projections $\Pi$ with $\rank\Pi\le cI$,
\begin{equation}
 \E|\Pi w|^2=o(I).
 \label{eq:prediction-small-subspaces}
\end{equation}
\item For every $\epsilon>0$, for all sufficiently far members of the
sequence there is a deterministic subspace $V_\epsilon$ of dimension
at most $C_\epsilon I$ such that, for every $v\perp V_\epsilon$,
\begin{equation}
 v^*(\E ww^*)v\le\epsilon|v|^2,\qquad
 v^*(\E b_sb_s^*)v\le\epsilon|v|^2,\qquad
 v^*(\E M_s)v\le\epsilon|v|^2.
 \label{eq:prediction-outlier-subspaces}
\end{equation}
In the gap case $V_\epsilon$ may be taken to be zero, and
$\|\E b_sb_s^*\|_{\op}+\|\E M_s\|_{\op}=o(1)$.
\item In the strict case, for every fixed $r\in(0,1/4)$ and
$\epsilon>0$, the conditional covariance
$C_r=\Cov(w(Y_0)\mid U_r)$ satisfies
\begin{equation}
 \E\tr\bigl(C_r\1_{\{C_r>\epsilon\Id\}}\bigr)=o(I).
 \label{eq:conditional-tail}
\end{equation}
The indicator denotes the spectral projection.  No assertion about
additional Gaussian probes or Equation~\eqref{eq:conditional-tail}
is needed for the gap case.
\end{enumerate}
\end{lemma}

\begin{proof}
We first prove the strict-case spectral estimates.  If a random vector
$S$ has covariance $C$ and mean $a$, the affine estimator
\[
 a+C(\Id+C)^{-1}(S+G-a)
\]
of $S$ from $S+G$ is its least-squares affine predictor.  Its mean
squared norm is $|a|^2+\tr(C^2(\Id+C)^{-1})$, by direct expansion of the
covariances.  Conditional expectation has at least this squared
norm, because it is the orthogonal projection onto all functions of
the observation.  The same calculation holds conditionally on any
side information, using fresh independent Gaussian noise.

By the Gaussian signal derivative formula and its monotonicity,
\[
 \E|\E[w\mid L_1]|^2
 \le2\bigl(d(2)-d(1)\bigr)\le2d(2)=o(I).
\]
Apply the affine prediction bound to $S=w$ and $C=\Cov_\eta w$.
For every fixed $\epsilon>0$,
\begin{equation}
 \tr\bigl(C\1_{\{C>\epsilon\Id\}}\bigr)=o(I),
 \label{eq:prediction-unconditional-tail}
\end{equation}
since $t^2/(1+t)\ge\epsilon t/(1+\epsilon)$ for $t>\epsilon$.
For every projection of rank at most $cI$,
\[
 \E|\Pi w|^2
 \le\epsilon cI+\tr\bigl(C\1_{\{C>\epsilon\Id\}}\bigr)
            +|\E w|^2.
\]
First take the sequence limit, then $\epsilon\downarrow0$.
Equation~\eqref{eq:prediction-input} proves
Equation~\eqref{eq:prediction-small-subspaces} uniformly in $\Pi$.
In the gap case a projection of rank at most $cI$ is eventually zero,
which proves the same assertion.

Conditionally on $U_r$, the additional observation $L_1$ has signal
$w$ and covariance $C_r$.  Orthogonality of successive conditional
expectations and the conditional affine prediction bound yield
\[
 I\bigl(p_1(r)-m(r)\bigr)
 \ge\E\tr\bigl(C_r^2(\Id+C_r)^{-1}\bigr).
\]
Proposition~\ref{prop:prediction} makes the left side $o(I)$.
The same scalar eigenvalue inequality proves
Equation~\eqref{eq:conditional-tail}.

We finally bound the dimensions of the deterministic outlier spaces.
Let
\[
 D=\E_\eta\E[XX^*\mid X+\sqrt sG=Y_0],
\]
where the inner posterior kernel always refers to the original prior
$\mu$.  Conditional Jensen and the definition of $M_s$ imply
\begin{equation}
 \E_\eta b_sb_s^*\preceq D,\qquad
 \E_\eta M_s\preceq D/s.
 \label{eq:prediction-posterior-domination}
\end{equation}
For each deterministic projection $Q$ of rank $l\ge1$, the original
subgaussian assumption gives
\begin{equation}
 \log\E_\mu\exp\left(\frac{|QX|^2}{Ca^2}\right)\le Cl.
 \label{eq:prediction-subspace-moment}
\end{equation}
Indeed choose a $1/2$-net of the unit sphere in
$\operatorname{ran}Q$ with at most $5^l$ points.  The norm $|QX|$ is
at most twice the largest absolute linear form over this net, and
bounding the exponential of a maximum by the sum proves the claim.

In the strict case, form the joint law by first drawing $Y_0$ from
$\eta$ and then drawing $X$ from the original posterior kernel.
Its entropy relative to the original joint law of
$(X,X+\sqrt sG)$ is exactly $H_0=I/(2\lambda)$.  Data processing
therefore bounds the entropy of its $X$-marginal relative to $\mu$ by
$H_0$.  Applying the entropy variational inequality to
Equation~\eqref{eq:prediction-subspace-moment} gives
\begin{equation}
 \tr(QD)\le Ca^2(H_0+l).
 \label{eq:prediction-posterior-trace}
\end{equation}
If the spectral subspace of $D$ above $\tau>0$ has dimension $l$,
then $\tau l\le Ca^2(H_0+l)$.  For sufficiently small $a$ this implies
$l\le2Ca^2H_0/\tau=O_\tau(I)$; the assertion is trivial when $l=0$.
The spectral subspace of $\E ww^*$ above $\epsilon$ has dimension
at most $I/\epsilon$, since its trace is $I$.  Take $V_\epsilon$ to
be the span of this subspace and the spectral subspace of $D$ above
$s\epsilon$.  Its dimension is at most $C_\epsilon I$, and
Equation~\eqref{eq:prediction-posterior-domination} gives all three
bounds in Equation~\eqref{eq:prediction-outlier-subspaces}.

In the gap case, averaging the original posterior under
$\eta=\mu_s$ gives $D=\E_\mu XX^*\preceq Ca^2\Id$.  Thus
Equation~\eqref{eq:prediction-posterior-domination} gives the two
claimed operator bounds directly.  Also
$\|\E ww^*\|_{\op}\le I\to0$, so $V_\epsilon=\{0\}$ suffices
eventually.  This completes all assertions.
\end{proof}
